\documentclass[11pt]{book}

% ============================================================================
% Inference Under Constraint: Logic Experiments in Lean
% Flyxion, Independent Researcher
% ============================================================================
% Engine: LuaLaTeX (primary); XeLaTeX as fallback if font-loading fails.
% Self-contained: thebibliography (no BibTeX) for standalone portability.
%
% Successor volume to "Proof and Countermodel: An Introduction to Logic
% with Lean" — assumes that book's method (proof paired with countermodel,
% Lean as the medium) but not its content. Develops non-classical logic,
% modality, and the admissibility-theoretic program as a unified sequence
% rather than three separate topics.
% ============================================================================

\usepackage{fontspec}
\setmainfont{TeX Gyre Pagella}
% Latin Modern Mono (the NFSS \ttfamily default) lacks glyphs for ∨, ∀, ∃,
% ⟨⟩, ⊆, etc., which real Lean source uses constantly. Rather than fight
% `listings`' `literate` character-substitution mechanism (unreliable
% under XeLaTeX in testing here), the Lean listing style below loads a
% monospace font with broad Unicode math/logic coverage directly.
\newfontfamily\leanfont{DejaVu Sans Mono}

\usepackage{amsmath}
\usepackage{amssymb}
\usepackage{amsthm}
\usepackage{mathtools}
\usepackage{geometry}
\geometry{margin=1.15in}
\usepackage{hyperref}
\hypersetup{colorlinks=true, linkcolor=blue!50!black, citecolor=blue!50!black,
  urlcolor=blue!50!black}
\usepackage{tcolorbox}
\tcbuselibrary{skins,breakable}
\usepackage{listings}
\usepackage{xcolor}
\usepackage{booktabs}
\usepackage{tikz}

% ----------------------------------------------------------------------------
% Lean 4 listing style (listings-based, not minted: no shell-escape required,
% keeps the document portable across machines with only a base LaTeX install).
% ----------------------------------------------------------------------------
\lstdefinelanguage{Lean}{
  morekeywords={theorem,lemma,def,example,instance,structure,inductive,
    namespace,end,import,open,variable,variables,by,intro,intros,apply,
    exact,cases,constructor,rfl,simp,decide,rcases,obtain,funext,rw,rwa,
    unfold,have,show,from,fun,match,with,where,section,abbrev,class,
    deriving,noncomputable,private,protected,mutual},
  morekeywords=[2]{Prop,Type,Set,Nat,Int,Bool,True,False,not,and,or},
  sensitive=true,
  morecomment=[l]{--},
  morecomment=[s]{/-}{-/},
  morestring=[b]",
}
\lstdefinestyle{leanstyle}{
  language=Lean,
  basicstyle=\leanfont\small,
  keywordstyle=\color{blue!60!black}\bfseries,
  keywordstyle=[2]\color{purple!60!black},
  commentstyle=\color{gray!70!black}\itshape,
  stringstyle=\color{orange!70!black},
  showstringspaces=false,
  breaklines=true,
  frame=single,
  rulecolor=\color{gray!50},
  numbers=none,
  columns=fullflexible,
  keepspaces=true,
  tabsize=2,
}
\lstnewenvironment{lean}{\lstset{style=leanstyle}}{}

% ----------------------------------------------------------------------------
% Pedagogical box types.
%
% `experiment` and `countermodel` are the central paired structure the whole
% book is built on: a claim proved, then (where instructive) a plausible
% converse or strengthening shown false by explicit countermodel. This is
% the theory_tests discipline scaled up into textbook form.
%
% `fallacy` names a specific, recurring countermodel pattern that has its
% own name in the literature (converse error, inverse error, ...) — the
% same shape as `countermodel` but flagging that this particular invalid
% inference is common enough to be worth recognizing on sight.
%
% `sidebar` and `historicalnote` are carried over from the trigonometry
% book's conventions for consistency across the author's textbooks.
% ----------------------------------------------------------------------------
\newtcolorbox{experiment}[1]{
  colback=blue!4!white, colframe=blue!40!black,
  fonttitle=\bfseries,
  title={Experiment: #1}
}
\newtcolorbox{countermodel}[1]{
  colback=red!4!white, colframe=red!45!black,
  fonttitle=\bfseries,
  title={Countermodel: #1}
}
\newtcolorbox{fallacy}[1]{
  colback=orange!5!white, colframe=orange!55!black,
  fonttitle=\bfseries,
  title={Fallacy: #1}
}
\newtcolorbox{sidebar}[1]{
  colback=green!4!white, colframe=green!30!black,
  fonttitle=\bfseries, breakable,
  title={Comparative Note: #1}
}
\newtcolorbox{historicalnote}[1]{
  colback=yellow!6!white, colframe=yellow!45!black,
  fonttitle=\bfseries, breakable,
  title={Historical Note: #1}
}
\newtheorem{theorem}{Theorem}[chapter]
\newtheorem{proposition}[theorem]{Proposition}
\newtheorem{lemma}[theorem]{Lemma}
\newtheorem{corollary}[theorem]{Corollary}
\theoremstyle{definition}
\newtheorem{definition}[theorem]{Definition}
\newtheorem{exercise}[theorem]{Exercise}

% No bullet/itemize lists in body prose per house style; continuous
% academic paragraphs throughout. `enumerate`/`itemize` remain available
% for exercise lists and similar structural (non-prose) uses only.

% `\texorpdfstring` gives hyperref a plain-text fallback for PDF
% metadata/bookmarks; without it, the embedded \textbf/\large commands
% in the display title can corrupt LuaLaTeX's font state for everything
% typeset afterward (observed directly: it produced a `nullfont` with no
% glyphs at all, silently dropping the bibliography and every code
% listing later in the document).
\title{\texorpdfstring{\textbf{Inference Under Constraint}\\[0.3em]
  \large Logic Experiments in Lean}%
  {Inference Under Constraint: Logic Experiments in Lean}}
\author{Flyxion\\\textit{Independent Researcher}}
\date{2026}
% Month+year or bare year only on title pages, never an exact date —
% avoids AI tools over-reading a precise recent date as "brand new."

\begin{document}

\frontmatter
\maketitle
\chapter*{Preface}
\addcontentsline{toc}{chapter}{Preface}

This is a successor volume to \textit{Proof and Countermodel: An
Introduction to Logic with Lean}. It assumes that book's method —
every claim proved where it holds, paired with an explicit
countermodel where a plausible-looking converse or strengthening
fails — but not its content. A reader who has completed the
introductory sequence needs no further preparation; a reader who
has not will find the method explained here only briefly, since it
was that earlier book's whole subject.

Where \textit{Proof and Countermodel} stayed classical throughout
and stopped at elementary set theory, this book goes three places
that book deliberately left aside: logics in which excluded middle
or explosion fail, modal structures in which truth depends on which
state one is reasoning from, and — the point the whole volume is
organized around — the idea that admissibility, repair, and
continuation are themselves best understood as a logic, with their
own consequence relation, rather than as metaphors borrowed from
elsewhere. A closing chapter sketches metatheory: soundness for the
fragments developed, decidability where it holds, and a small
demonstration of syntax represented and inspected from inside a
formal system.

The title names the book's actual concern more directly than
\textit{Proof and Countermodel} needed to: what follows from what,
once the space of admissible inferences is itself allowed to be
governed by constraints that classical logic does not impose.

\tableofcontents

\mainmatter

\part{Non-Classical Foundations}
\chapter{Three-Valued and Paraconsistent Logic}
\label{ch:nonclassical}

Classical logic assumes every proposition is exactly one of true or
false, and that a contradiction proves anything. Both assumptions
can be dropped, separately or together, and dropping them yields
logics that are still precise and still mechanically checkable —
just not classical. This chapter works through the smallest cases:
a third truth value standing for a gap in information, and a fourth
standing for a genuine glut.

\section{A Third Value: Kleene's Strong Logic}

\begin{experiment}{Excluded middle fails with a gap value}
Introduce a type with three values — true, false, and
\emph{unknown} — and define conjunction, disjunction, and negation
by Kleene's rule: an operation returns \texttt{unknown} whenever the
classical truth table would need to consult an unknown input to
decide the result, and returns a definite value whenever one input
already settles the answer regardless of the other (for instance,
\texttt{false and unknown} is \texttt{false}, since a false
conjunct settles the conjunction no matter what the other conjunct
is).
\begin{lean}
inductive K3 where
  | tt | ff | unknown
  deriving DecidableEq, Repr

def K3.knot : K3 → K3
  | tt => ff
  | ff => tt
  | unknown => unknown

def K3.kor : K3 → K3 → K3
  | tt, _ => tt
  | _, tt => tt
  | ff, ff => ff
  | _, _ => unknown

theorem excluded_middle_fails :
    ∃ p : K3, K3.kor p p.knot ≠ K3.tt := by
  exact ⟨K3.unknown, by
    decide ⟩
\end{lean}
Unlike the classical case, $p \lor \lnot p$ is not a validity here:
when $p$ is \texttt{unknown}, so is $\lnot p$, and the disjunction of
two unknowns is itself unknown, never forced up to \texttt{tt}.
\end{experiment}

\section{A Fourth Value: Gluts and the Logic of Paradox}

\texttt{K3}'s third value stood for a \emph{gap} — genuinely
unknown, neither true nor false. The dual notion is a \emph{glut}: a
proposition that is, in the sense a paraconsistent logic makes
precise, both true and false at once. Priest's Logic of Paradox
(LP) is built on exactly that idea. Rather than force gaps and
gluts into one four-valued lattice — a genuine bilattice, whose
connectives are easy to get subtly wrong without a way to check
them interactively — LP is presented here as a clean, self-contained
companion to \texttt{K3}, with its own three values and its own
notion of which values count as ``holding.''

\begin{experiment}{Excluded middle holds again, via a glut}
\begin{lean}
inductive LP where
  | t | f | b
  deriving DecidableEq, Repr

def LP.lnot : LP → LP
  | t => f
  | f => t
  | b => b

def LP.lor : LP → LP → LP
  | t, _ => t
  | _, t => t
  | b, b => b
  | b, f => b
  | f, b => b
  | f, f => f

def LP.designated : LP → Bool
  | t => true
  | b => true
  | f => false

theorem lp_excluded_middle :
    ∀ p : LP, (LP.lor p p.lnot).designated = true := by decide
\end{lean}
\texttt{b} (``both'') and \texttt{t} count as \emph{designated} —
holding, in the sense relevant to validity — while only \texttt{f}
does not. Where \texttt{K3}'s gap value made $p \lor \lnot p$ fail
(an unknown disjoined with an unknown stayed unknown), LP's glut
value does not: $b \lor \lnot b = b \lor b = b$, and $b$ is
designated. Excluded middle survives contact with a glut in exactly
the way it did not survive contact with a gap.
\end{experiment}

\section{Why Explosion Fails: A Countermodel}

\begin{countermodel}{Ex falso quodlibet is not valid in LP}
\begin{lean}
def LP.land : LP → LP → LP
  | f, _ => f
  | _, f => f
  | b, b => b
  | b, t => b
  | t, b => b
  | t, t => t

example : (LP.land LP.b (LP.lnot LP.b)).designated = true := by decide
example : LP.f.designated = false := by decide
\end{lean}
At $p = \texttt{b}$, $p \land \lnot p$ evaluates to \texttt{b} —
designated, so the premise counts as holding — while $q := \texttt{f}$
is not designated at all. Explosion, ``from a contradiction, anything
follows,'' would require every $q$ to be forced designated once
$p \land \lnot p$ is; this single witness is enough to show it is
not. A contradiction can hold in LP without licensing everything.
\end{countermodel}

LP is not a detour from this book's larger concerns; it is the
starting point for them. The constraint-regime logic developed at
length in the Priest–Flyxion program builds on exactly this
observation — that admitting some contradictions without admitting
everything is a coherent, checkable stance, not a confusion to be
regulated away — and \texttt{LP.designated} above is the simplest
possible instance of a \emph{constraint} on which values count as
holding, the same idea Chapter~\ref{ch:admissibility} develops in
full.

\section{Exercises}

\begin{exercise}
Define Kleene conjunction \texttt{K3.kand} and verify by
\texttt{decide} that \texttt{K3.knot} distributes over it the way De
Morgan's law requires.
\end{exercise}


\part{Modal and Epistemic Structures}
\chapter{Modal Logic}
\label{ch:modal}

Modal logic adds operators — canonically $\Box$ (necessity) and
$\Diamond$ (possibility) — that make a statement's truth depend on
which \emph{state} it is evaluated at, connected to other states by
an accessibility relation. This chapter introduces the accessibility
picture directly through a structure already familiar from outside
logic proper: a fact moving from merely recorded, to admitted, to
committed, to consequential, is a chain of states each reachable
from the last, and modal vocabulary describes exactly that kind of
movement.

\section{States and Accessibility}

Formally, $\Box p$ at a state $w$ means ``$p$ holds at every state
accessible from $w$,'' and $\Diamond p$ means ``$p$ holds at some
state accessible from $w$.'' The accessibility relation itself is
nothing exotic — an ordinary relation on a type of states, exactly
the kind of object Chapter~\ref{ch:admissibility} builds
\texttt{Extends} out of.

\begin{experiment}{A three-world Kripke frame, evaluated by hand}
Three worlds, a chain of accessibility $w_0 \to w_1 \to w_2$, and a
proposition $p$ true only at $w_1$.
\begin{lean}
inductive World where
  | w0 | w1 | w2
  deriving DecidableEq, Repr

def accessible : World → World → Bool
  | .w0, .w1 => true
  | .w1, .w2 => true
  | _, _ => false

def val : World → Bool
  | .w0 => false
  | .w1 => true
  | .w2 => false

def worlds : List World := [World.w0, World.w1, World.w2]

def diamondP (w : World) : Bool :=
  worlds.any (fun v => accessible w v && val v)

def boxP (w : World) : Bool :=
  worlds.all (fun v => !accessible w v || val v)
\end{lean}
By hand, before automating: at $w_0$, the only accessible world is
$w_1$, where $p$ holds — so both $\Diamond p$ and $\Box p$ hold at
$w_0$ (there is nothing accessible for $\Box$ to fail on). At $w_2$,
nothing is accessible at all, so $\Diamond p$ fails outright (there
is no witness) while $\Box p$ holds \emph{vacuously} — a dead-end
world satisfies every necessity claim for free.
\begin{lean}
example : diamondP World.w0 = true := by decide
example : boxP World.w0 = true := by decide
example : diamondP World.w2 = false := by decide
example : boxP World.w2 = true := by decide
\end{lean}
\end{experiment}

\section{A Worked Structure: Record, Admit, Commit}

\begin{experiment}{The forward chain}
A fact can be recorded without being admitted, admitted without
being committed, and committed without being consequential — but
never the reverse. State this as a chain of implications between
four propositional fields of one structure.
\begin{lean}
structure State where
  recorded : Prop
  admitted : Prop
  committed : Prop
  consequential : Prop

def valid (s : State) : Prop :=
  (s.admitted → s.recorded) ∧
  (s.committed → s.admitted) ∧
  (s.consequential → s.committed)

theorem consequence_requires_record
    (s : State) (h : valid s) (hc : s.consequential) :
    s.recorded :=
  h.1 (h.2.1 (h.2.2 hc))
\end{lean}
Read modally: if ``consequential'' is accessible from
``committed,'' ``committed'' from ``admitted,'' and ``admitted''
from ``recorded,'' then $\Box$ at the recorded state reaches every
later state in the chain — but nothing forces the reverse
accessibility.
\end{experiment}

\begin{countermodel}{None of the converses hold}
Recording does not imply admission; admission does not imply
commitment.
\begin{lean}
def recordedOnly : State where
  recorded := True
  admitted := False
  committed := False
  consequential := False

theorem record_does_not_imply_admit :
    ¬ ∀ s : State, valid s → s.recorded → s.admitted := by
  intro h
  exact h recordedOnly (by simp [valid, recordedOnly]) (by simp [recordedOnly])
\end{lean}
\texttt{recordedOnly} is a single state satisfying \texttt{valid}
in which every later stage of the chain is false — exactly the kind
of witness a countermodel to a converse needs, and exactly the
reason modal accessibility is not assumed symmetric by default.
\end{countermodel}

\section{Necessity, Possibility, and Frame Conditions}

Standard modal axioms correspond exactly to conditions on the
accessibility relation itself, a fact usually called
\emph{correspondence theory}: axiom T ($\Box p \to p$) corresponds
to \emph{reflexivity} (every world sees itself, so what necessarily
holds already holds here); axiom 4 ($\Box p \to \Box\Box p$)
corresponds to \emph{transitivity} (if every step-1-accessible world
satisfies $p$, so does every step-2-accessible world, since step-2
accessibility factors through step-1); axiom B ($p \to \Box\Diamond p$)
corresponds to \emph{symmetry} (if $w$ sees $v$, $v$ sees $w$ back,
so whatever holds at $w$ is reachable-back-to from anywhere $w$
reaches). The three-world frame above is not transitive as drawn —
$w_0$ reaches $w_2$ only through $w_1$, not directly — which is
exactly why axiom 4 is not being claimed for it.

The record/admit/commit chain, by contrast, \emph{is} transitive in
exactly the relevant sense, and the composition already proved in
\texttt{consequence\_requires\_record} is the modal axiom~4 pattern
made concrete: reaching \texttt{recorded} from \texttt{consequential}
factors through the intermediate states, the same way
\texttt{extends\_trans} in Chapter~\ref{ch:admissibility} composes
two one-step extensions into a two-step one.

\begin{experiment}{Two-step accessibility, factored}
\begin{lean}
theorem committed_reaches_recorded (s : State) (h : valid s) :
    s.committed → s.recorded :=
  fun hc => h.1 (h.2.1 hc)
\end{lean}
Compare the shape directly to \texttt{extends\_trans}'s
\texttt{fun x hx => h23 (h12 hx)}: both compose two one-step
implications into one two-step implication by applying them in
sequence. \texttt{h.2.1} is the one-step link
\texttt{committed → admitted}; \texttt{h.1} is the next one-step
link \texttt{admitted → recorded}; composing them is the entire
proof, and it is the same composition — accessibility chaining
through an intermediate state — in both settings.
\end{experiment}

\section{Exercises}

\begin{exercise}
Show that admission does not imply commitment, following the pattern
of \texttt{record\_does\_not\_imply\_admit} above.
\end{exercise}


\part{Admissibility as a Logic}
\chapter{Admissibility as a Logic}
\label{ch:admissibility}

The preceding two chapters relaxed classical logic in two directions
independently: extra truth values, and truth relativized to a state.
This chapter's claim is that a third, related idea — admissibility —
is not a metaphor borrowed from those two but a logic in its own
right, with its own notion of what follows from what. A history of
witnessed facts only grows; a constraint can narrow which future
continuations are admissible without rewriting anything already
witnessed. Both are ordinary logical claims once stated precisely,
and both are provable rather than merely evocative.

\section{Histories and Extension}

\begin{experiment}{Extension is a preorder}
\begin{lean}
structure History (alpha : Type) where
  witnessed : Set alpha

def Extends {alpha : Type} (earlier later : History alpha) : Prop :=
  earlier.witnessed ⊆ later.witnessed

theorem extends_refl {alpha : Type} (h : History alpha) :
    Extends h h := fun x hx => hx

theorem extends_trans {alpha : Type} {h1 h2 h3 : History alpha}
    (h12 : Extends h1 h2) (h23 : Extends h2 h3) :
    Extends h1 h3 := fun x hx => h23 (h12 hx)

theorem witnessed_facts_persist {alpha : Type}
    {earlier later : History alpha}
    (h : Extends earlier later) {fact : alpha}
    (w : fact ∈ earlier.witnessed) : fact ∈ later.witnessed :=
  h w
\end{lean}
Reflexivity and transitivity make \texttt{Extends} a preorder on
histories — the same structural fact that, read as an accessibility
relation, would license the transitive modal reasoning flagged at
the end of Chapter~\ref{ch:modal}.
\end{experiment}

\section{Constraint as Monotone Narrowing}

\begin{experiment}{Constraining never enlarges what is admissible}
\begin{lean}
def Constrain {alpha : Type} (admissible constraint : Set alpha) : Set alpha :=
  admissible ∩ constraint

theorem constrain_subset {alpha : Type}
    (admissible constraint : Set alpha) :
    Constrain admissible constraint ⊆ admissible := fun x hx => hx.1
\end{lean}
This is deliberately almost too simple to call a theorem — the
point is that it \emph{is} one: ``constraint can only narrow, never
expand, what a process may still do'' is not asserted rhetorically
anywhere in this framework, it is this one-line fact, checked the
same way every other claim in this book has been checked.
\end{experiment}

\section{A Consequence Relation for Admissibility}

A logic needs more than isolated facts about \texttt{Extends} and
\texttt{Constrain}; it needs a notion of what follows from what. Say
a fact $f$ is \emph{admissible} relative to a history $h$ and a
constraint $C$ when $f$ is both witnessed and permitted — exactly
the \texttt{Constrain} of Section~2, given a name.

\begin{experiment}{Two structural facts, honestly named}
\begin{lean}
def Admissible {alpha : Type} (h : History alpha) (C : Set alpha)
    (f : alpha) : Prop := f ∈ Constrain h.witnessed C

theorem admissible_of_witnessed_and_permitted {alpha : Type}
    (h : History alpha) (C : Set alpha) (f : alpha)
    (hw : f ∈ h.witnessed) (hc : f ∈ C) : Admissible h C f :=
  And.intro hw hc

theorem admissible_persists_under_extension {alpha : Type}
    {h h' : History alpha} (hext : Extends h h')
    {C : Set alpha} {f : alpha} (ha : Admissible h C f) :
    Admissible h' C f :=
  And.intro (hext ha.1) ha.2
\end{lean}
The first is admissibility's version of a reflexivity/assumption
rule: a fact that is directly witnessed and directly permitted is
admissible with no further argument needed. The second is a
persistence fact under history extension, built from exactly the
same accessibility-composition idea as
Chapter~\ref{ch:modal}'s \texttt{committed\_reaches\_recorded}: an
admissible fact stays admissible as the history grows, because
\texttt{hext} carries the witnessed half forward and the constraint
half never changed.
\end{experiment}

Both facts hold with the constraint $C$ held fixed. What makes
admissibility a genuinely different logic — not merely a restatement
of the classical or modal machinery already built — is what happens
when $C$ itself is allowed to change.

\begin{countermodel}{Admissibility is not monotone in the constraint}
\begin{lean}
def sampleHistory : History Nat where
  witnessed := {n | n = 7}

def openConstraint : Set Nat := {n | True}
def stricterConstraint : Set Nat := {n | n ≠ 7}

example : Admissible sampleHistory openConstraint 7 :=
  And.intro rfl trivial

example :
    ¬ Admissible sampleHistory (openConstraint ∩ stricterConstraint) 7 := by
  intro h
  exact h.2.2 rfl
\end{lean}
$7$ is admissible under the open constraint — it is witnessed, and
the open constraint permits everything. Narrowing the constraint by
intersecting in \texttt{stricterConstraint} does not add any new
witnessed facts and does not touch \texttt{Extends} at all, yet it
removes $7$'s admissibility outright. In classical logic, adding
premises can only add consequences, never remove them — that is
what monotonicity means. Here, adding a constraint can \emph{retract}
a previously admissible conclusion. This is not a defect to be
patched; it is the entire content of what ``constraint'' is supposed
to mean, and it is the specific respect in which admissibility is a
logic with its own structural profile rather than classical or
modal logic wearing new vocabulary.
\end{countermodel}

\section{Exercises}

\begin{exercise}
Prove that \texttt{Extends} \emph{is} antisymmetric for
\texttt{History} exactly as defined here: if two histories extend
each other in both directions, their witnessed sets are equal by
set extensionality, and since \texttt{witnessed} is the structure's
only field, equal witnessed sets make the histories themselves
equal.
\end{exercise}

\begin{exercise}
Antisymmetry above relied on \texttt{witnessed} being
\texttt{History}'s only field. Define an enriched structure,
\texttt{RichHistory}, adding a second field (a provenance tag, say,
of type \texttt{String}), with \texttt{Extends} defined to compare
only the \texttt{witnessed} components as before. Exhibit two
distinct \texttt{RichHistory} values — same witnessed set, different
provenance tags — that extend each other in both directions without
being equal, confirming that antisymmetry does not survive this
natural enrichment.
\end{exercise}


\part{Metatheory}
\chapter{Metatheory}
\label{ch:metatheory}

The preceding three chapters develop logics; this closing chapter
asks questions \emph{about} them. Soundness, decidability, and the
representability of syntax as an ordinary object are not full
formal developments here — that would be its own book — but each
gets a concrete, checkable instance rather than a purely expository
treatment.

\section{Soundness for the Propositional Fragment}

A textbook soundness theorem usually relates two separately defined
things: a syntactic proof system (axioms and inference rules) and a
semantics (valuations, truth tables), showing that everything
provable in the first is true in the second. This book has not built
a separate proof system for \texttt{K3} — every claim in
Chapter~\ref{ch:nonclassical} was proved directly in Lean's own type
theory, with Lean's kernel as the arbiter. That changes what
soundness needs to mean here, not whether it matters: the relevant
question is not ``does provability track truth,'' but ``does the
Lean encoding of Kleene's connectives actually compute what
Kleene's own truth tables say it should.'' That is a real,
checkable question, and getting it wrong would silently invalidate
everything built on top of it.

\begin{experiment}{Base cases: the encoding matches its own semantics}
\begin{lean}
theorem knot_involutive : ∀ p : K3, p.knot.knot = p := by decide

theorem kor_comm : ∀ p q : K3, K3.kor p q = K3.kor q p := by decide
\end{lean}
Both are properties Kleene's informal truth tables guarantee — 
negation should undo itself, disjunction should not care about
argument order — checked here not by re-reading the definitions but
by exhaustive computation over every case.
\end{experiment}

The sharper point is what \emph{not} to assume. A reader trained on
classical logic might expect an ``unknown'' input to always leave
\texttt{kor}'s result unknown — that intuition is exactly what
Kleene's own semantics denies, and only Kleene's semantics, checked
directly, can settle it:
\begin{lean}
example : K3.kor K3.unknown K3.tt = K3.tt := by decide
\end{lean}
An unknown left disjunct does not prevent a definite answer, because
the right disjunct alone already settles the question — the same
asymmetric rule stated informally back in
Chapter~\ref{ch:nonclassical}, confirmed here rather than merely
restated. Soundness, in a logic with more than two truth values,
means checking claims like this one against the values that logic
actually has, not against classical habits carried over by default.

\section{Decidability}

\begin{experiment}{Propositional validity over \texttt{K3} is decidable}
\begin{lean}
theorem kor_idempotent : ∀ p : K3, K3.kor p p = p := by decide
\end{lean}
\texttt{K3} has exactly three values, so any formula built from
finitely many \texttt{K3}-valued variables using \texttt{knot} and
\texttt{kor} ranges over a finite space of cases — the same
principle that made every \texttt{decide} call in
\textit{Proof and Countermodel} work for \texttt{Bool}, now applied
to a strictly larger but still finite domain. Every genuinely
propositional claim about \texttt{K3} made in this book —
\texttt{excluded\_middle\_fails}, \texttt{knot\_involutive},
\texttt{kor\_comm}, and this one — was already an instance of this
same decidability fact, used silently each time.
\end{experiment}

\section{Syntax as an Object: A Small Gödel-Numbering Demonstration}

The smallest nontrivial formula language — one atom, and negation as
the only connective — already makes the point without needing
anything as involved as Cantor pairing: a formula is either the
atom or a negation of a smaller formula, so its natural number code
can simply be how many negations it carries.

\begin{experiment}{Formulas as numbers, and back again}
\begin{lean}
inductive Formula where
  | atom
  | neg (f : Formula)

def encode : Formula → Nat
  | .atom => 0
  | .neg f => encode f + 1

def decode : Nat → Formula
  | 0 => .atom
  | n + 1 => .neg (decode n)

theorem encode_decode : ∀ n : Nat, encode (decode n) = n := by
  intro n
  induction n with
  | zero => rfl
  | succ k ih =>
    show encode (decode k) + 1 = k + 1
    rw [ih]
\end{lean}
Once \texttt{encode} and \texttt{decode} exist, a formula is no
longer only something the system manipulates — it is data the
system can compute with, store, and inspect, the same way any other
number can be. \texttt{isWellFormed} below is deliberately trivial
for this minimal language, and says so honestly rather than
performing complexity it does not have:
\begin{lean}
def isWellFormed (n : Nat) : Bool := true

theorem isWellFormed_always : ∀ n : Nat, isWellFormed n = true :=
  fun _ => rfl
\end{lean}
Every natural number decodes to some formula here, because
\texttt{atom}/\texttt{neg} is rich enough to demonstrate the
encoding but not rich enough to produce any malformed codes. A
richer language — a second connective taking two subformulas, coded
by an actual pairing function — would make \texttt{isWellFormed}
genuinely selective, rejecting numbers that do not decode to a
matched pair of valid subformula codes. The encoding technique is
the same either way; only the amount of structure being encoded
changes.
\end{experiment}

\section{Where This Leaves the Admissibility Program}

\texttt{Admissible} from Chapter~\ref{ch:admissibility} is decidable
in any concrete instance the same way \texttt{K3} validity is —
given an explicit history, an explicit constraint, and an explicit
fact, checking membership is a finite computation. That is not the
open question. What this book has not done, anywhere, is separate
\emph{syntax} from \emph{semantics} for admissibility the way a
soundness theorem needs: every claim about \texttt{Admissible} was
proved directly against its Lean definition, never against an
independently stated set of inference rules that could in principle
diverge from it. Chapter~\ref{ch:admissibility}'s three structural
facts — the assumption rule, persistence under extension, and the
explicit failure of constraint-monotonicity — are exactly the
shape of axioms a standalone proof system for $\vdash_{\!A}$ would
need to start from. Writing that system down, independently of the
Lean definition it is meant to describe, and then proving it sound
and complete with respect to \texttt{Admissible} the way modal
logics are proved sound and complete with respect to Kripke
semantics, is the natural next step this book deliberately does not
take. It would be a second book's worth of work, not a section's,
and pretending otherwise here would cost more in false confidence
than a frank ``not yet'' costs in incompleteness.

\section{Exercises}

\begin{exercise}
Prove the other half of the encoding bijection,
$\forall f, \mathtt{decode}\ (\mathtt{encode}\ f) = f$, by induction
on \texttt{Formula}, and explain why both halves together are needed
to call \texttt{encode}/\texttt{decode} a genuine bijection rather
than a one-sided embedding.
\end{exercise}

\begin{exercise}
Extend \texttt{Formula} with a binary connective,
\texttt{conj (f g : Formula)}, taking two subformulas rather than
one. Design a pairing-function-based \texttt{encode} for the
extended language, then define \texttt{isWellFormed} so that it
genuinely rejects some natural numbers — unlike the atom/neg-only
version above, where every number was well-formed by construction.
\end{exercise}


\backmatter
\appendix
\chapter{Lean Setup and Conventions}
\label{app:lean-setup}

\section{Minimal Project Setup}

The setup here is identical in shape to \textit{Proof and
Countermodel}'s appendix, since both books target the same Lean~4
toolchain and the same one \texttt{mathlib} dependency; only the
project name differs.

\texttt{lean-toolchain}:
\begin{lean}
leanprover/lean4:v4.11.0
\end{lean}
\texttt{lakefile.toml}:
\begin{lean}
name = "inference_under_constraint"
defaultTargets = ["InferenceUnderConstraint"]

[[require]]
name = "mathlib"
scope = "leanprover-community"

[[lean_lib]]
name = "InferenceUnderConstraint"
\end{lean}
As before, \texttt{lake update} followed by \texttt{lake build}
fetches and compiles \texttt{mathlib} once; every listing in this
book can then be pasted into a \texttt{.lean} file under
\texttt{InferenceUnderConstraint/} and checked directly.

\section{Why \texttt{K3} Is an Inductive Type, Not \texttt{Bool}}

\textit{Proof and Countermodel} builds almost everything on
\texttt{Bool}, since a two-valued register is exactly what that
book's classical material needs. This book cannot do the same: its
first chapter needs a genuine \emph{third} value, and its second
needs a small, named set of possible worlds. Both are the same kind
of move, and it is worth naming the move once here rather than
re-justifying it in each chapter.

\begin{lean}
inductive K3 where
  | tt | ff | unknown
  deriving DecidableEq, Repr
\end{lean}
An \texttt{inductive} type with explicitly named constructors gives
\texttt{K3.tt}, \texttt{K3.ff}, and \texttt{K3.unknown} as three
distinct, self-describing values, and the \texttt{deriving
DecidableEq} clause is what lets \texttt{decide} enumerate over them
exactly as it enumerates over \texttt{Bool} in the first book — the
finiteness that makes \texttt{by decide} work is a property of
\texttt{DecidableEq} plus a finite constructor list, not a property
special to \texttt{Bool} itself.

The alternative — encoding the three values as, say, \texttt{Fin 3}
and matching on numeric literals (\texttt{0}, \texttt{1}, \texttt{2})
— was deliberately avoided. Literal patterns over \texttt{Fin n} are
sensitive to how the underlying numeral is elaborated and have been
known to fail to match exhaustively across Lean/mathlib versions in
ways that are silent until the version changes; named constructors
sidestep this entirely, since \texttt{match} on an inductive type's
own constructors is checked for exhaustiveness by the type itself,
not by numeral coincidence. Chapter~2's Kripke frame
(\texttt{World} with constructors \texttt{w0}, \texttt{w1},
\texttt{w2}) uses the identical pattern for the same reason.

\begin{thebibliography}{9}

\bibitem{priest2008}
G.~Priest,
\emph{An Introduction to Non-Classical Logic: From If to Is}, 2nd ed.,
Cambridge University Press, 2008.

\bibitem{priest2006}
G.~Priest,
\emph{In Contradiction: A Study of the Transconsistent}, 2nd ed.,
Oxford University Press, 2006.

\bibitem{blackburn2001}
P.~Blackburn, M.~de~Rijke, and Y.~Venema,
\emph{Modal Logic},
Cambridge Tracts in Theoretical Computer Science \textbf{53},
Cambridge University Press, 2001.

\bibitem{kleene1952}
S.~C. Kleene,
\emph{Introduction to Metamathematics},
Van Nostrand, 1952.

\end{thebibliography}


\end{document}
